\chapter{Boolean Analysis --- Threshold Functions --- Linear Threshold Influence}
%%% generated by the blueprint pipeline from TCSlib.BooleanAnalysis.ThresholdFunctions.LinearThresholdInfluence
\section{Overview}
Discrete differentiation identifies the influence of a unate sign function with its singleton Fourier coefficient. Affine threshold representations then give a square-root bound on total influence.

\section{Declarations}

\begin{lemma}[A coordinate update leaves or flips a bit]
\lean{BooleanAnalysis.ThresholdFunctions.update_eq_ite_flipBit}
\leanok
\uses{BooleanAnalysis.BoolCube, BooleanAnalysis.flipBit}
\difficulty{4}
Let $n\ge0$, $x\in\{0,1\}^n$, $i\in[n]$, and $b\in\{0,1\}$. Replacing $x_i$ by $b$ leaves $x$ unchanged if $x_i=b$ and otherwise gives the vector obtained by flipping bit $i$.
\end{lemma}

\begin{lemma}[Discrete derivative through a bit flip]
\lean{BooleanAnalysis.ThresholdFunctions.derivative_eq_flipBit}
\leanok
\uses{BooleanAnalysis.BoolCube, BooleanAnalysis.BooleanFunc, BooleanAnalysis.ThresholdFunctions.update_eq_ite_flipBit, BooleanAnalysis.boolToSign, BooleanAnalysis.derivative, BooleanAnalysis.flipBit}
\difficulty{4}
Let $n\ge0$ and give $\{0,1\}^n$ uniform probability; write $\mathbb E$ for expectation. For a real cube function $h$, write $D_i h(x)=\tfrac12(h(x_{i\leftarrow0})-h(x_{i\leftarrow1}))$, where $x_{i\leftarrow b}$ replaces bit $i$ by $b$. For every coordinate $i$ and point $x$, $D_i h(x)=\tfrac12(h(x)-h(x^{\{i\}}))(-1)^{x_i}$, where $x^{\{i\}}$ flips bit $i$.
\end{lemma}

\begin{lemma}[Influence is the derivative second moment]
\lean{BooleanAnalysis.ThresholdFunctions.influence_eq_expect_derivative_sq}
\leanok
\uses{BooleanAnalysis.BooleanFunc, BooleanAnalysis.ThresholdFunctions.derivative_eq_flipBit, BooleanAnalysis.boolToSign_sq, BooleanAnalysis.derivative, BooleanAnalysis.expect, BooleanAnalysis.influence}
\difficulty{3}
Let $n\ge0$ and give $\{0,1\}^n$ uniform probability; write $\mathbb E$ for expectation. For a real cube function $h$, write $D_i h(x)=\tfrac12(h(x_{i\leftarrow0})-h(x_{i\leftarrow1}))$, where $x_{i\leftarrow b}$ replaces bit $i$ by $b$. Set $I_i(h)=\tfrac14\mathbb E[(h(x)-h(x^{\{i\}}))^2]$ and $I(h)=\sum_i I_i(h)$, where $x^{\{i\}}$ flips bit $i$. For every real cube function $h$ and $i\in[n]$, $I_i(h)=\mathbb E[(D_i h)^2]$.
\end{lemma}

\begin{lemma}[Coordinate averaging preserves the mean]
\lean{BooleanAnalysis.ThresholdFunctions.expect_expectationOperator}
\leanok
\uses{BooleanAnalysis.BoolCube, BooleanAnalysis.BooleanFunc, BooleanAnalysis.ThresholdFunctions.update_eq_ite_flipBit, BooleanAnalysis.expect, BooleanAnalysis.expectationOperator, BooleanAnalysis.flipBit, BooleanAnalysis.flipBit_flipBit}
\difficulty{4}
Let $n\ge0$ and give $\{0,1\}^n$ uniform probability; write $\mathbb E$ for expectation. Let $g$ be a real cube function and $i\in[n]$, and define $E_i g(x)=\tfrac12(g(x_{i\leftarrow0})+g(x_{i\leftarrow1}))$, where the indicated update replaces bit $i$. Then $\mathbb E[E_i g]=\mathbb E g$.
\end{lemma}

\begin{lemma}[Singleton Fourier coefficient as a derivative mean]
\lean{BooleanAnalysis.ThresholdFunctions.fourierCoeff_singleton_eq_expect_derivative}
\leanok
\uses{BooleanAnalysis.BooleanFunc, BooleanAnalysis.ThresholdFunctions.expect_expectationOperator, BooleanAnalysis.boolToSign, BooleanAnalysis.derivative, BooleanAnalysis.expect, BooleanAnalysis.expectationOperator, BooleanAnalysis.fourierCoeff, BooleanAnalysis.innerProduct}
\difficulty{3}
Let $n\ge0$ and give $\{0,1\}^n$ uniform probability; write $\mathbb E$ for expectation. Put $\chi_S(x)=\prod_{i\in S}(-1)^{x_i}$ and $\widehat h(S)=\mathbb E[h\chi_S]$ for $S\subseteq[n]$ and real cube functions $h$. For a real cube function $h$, write $D_i h(x)=\tfrac12(h(x_{i\leftarrow0})-h(x_{i\leftarrow1}))$, where $x_{i\leftarrow b}$ replaces bit $i$ by $b$. For every real cube function $h$ and coordinate $i$, $\widehat h(\{i\})=\mathbb E[D_i h]$.
\end{lemma}

\begin{lemma}[Square and absolute value of a sign derivative]
\lean{BooleanAnalysis.ThresholdFunctions.derivative_sq_eq_abs_of_pmOne}
\leanok
\uses{BooleanAnalysis.BoolCube, BooleanAnalysis.BooleanFunc, BooleanAnalysis.derivative, BooleanAnalysis.isPmOne}
\difficulty{4}
Let $n\ge0$ and give $\{0,1\}^n$ uniform probability; write $\mathbb E$ for expectation. For a real cube function $h$, write $D_i h(x)=\tfrac12(h(x_{i\leftarrow0})-h(x_{i\leftarrow1}))$, where $x_{i\leftarrow b}$ replaces bit $i$ by $b$. For any $h\colon\{0,1\}^n\to\{-1,1\}$, coordinate $i$, and cube point $x$, $(D_i h(x))^2=|D_i h(x)|$.
\end{lemma}

\begin{lemma}[Influence formula for a unate coordinate]
\lean{BooleanAnalysis.ThresholdFunctions.influence_eq_abs_fourierCoeff_of_sign}
\leanok
\uses{BooleanAnalysis.BooleanFunc, BooleanAnalysis.ThresholdFunctions.derivative_sq_eq_abs_of_pmOne, BooleanAnalysis.ThresholdFunctions.expect_neg, BooleanAnalysis.ThresholdFunctions.expect_nonneg, BooleanAnalysis.ThresholdFunctions.fourierCoeff_singleton_eq_expect_derivative, BooleanAnalysis.ThresholdFunctions.influence_eq_expect_derivative_sq, BooleanAnalysis.derivative, BooleanAnalysis.fourierCoeff, BooleanAnalysis.influence, BooleanAnalysis.isPmOne}
\difficulty{4}
Let $n\ge0$ and give $\{0,1\}^n$ uniform probability; write $\mathbb E$ for expectation. Put $\chi_S(x)=\prod_{i\in S}(-1)^{x_i}$ and $\widehat h(S)=\mathbb E[h\chi_S]$ for $S\subseteq[n]$ and real cube functions $h$. For a real cube function $h$, write $D_i h(x)=\tfrac12(h(x_{i\leftarrow0})-h(x_{i\leftarrow1}))$, where $x_{i\leftarrow b}$ replaces bit $i$ by $b$. Set $I_i(h)=\tfrac14\mathbb E[(h(x)-h(x^{\{i\}}))^2]$ and $I(h)=\sum_i I_i(h)$, where $x^{\{i\}}$ flips bit $i$. If $h$ is sign-valued and either $D_i h(x)\ge0$ for all $x$ or $D_i h(x)\le0$ for all $x$, then $I_i(h)=|\widehat h(\{i\})|$.
\end{lemma}

\begin{definition}[Walsh coefficients of an affine form]
\lean{BooleanAnalysis.ThresholdFunctions.affinePoly}
\leanok
\uses{BooleanAnalysis.ThresholdFunctions.MultilinearPolynomial}
For $n\ge0$, $c\in\mathbb R$, and $a\in\mathbb R^n$, the affine coefficient family $A_{c,a}$ has coefficient $c$ at $\varnothing$, coefficient $a_i$ at $\{i\}$, and zero at every set of size greater than one.
\end{definition}

\begin{lemma}[Constant coefficient of the affine family]
\lean{BooleanAnalysis.ThresholdFunctions.affinePoly_empty}
\leanok
\uses{BooleanAnalysis.ThresholdFunctions.affinePoly}
\difficulty{1}
Let $n\ge0$. Let $c\in\mathbb R$, $a\in\mathbb R^n$, and let $A_{c,a}(S)$ equal $c$ for $S=\varnothing$, $a_i$ for $S=\{i\}$, and zero for all larger $S$. Then $A_{c,a}(\varnothing)=c$.
\end{lemma}

\begin{lemma}[Singleton coefficient of the affine family]
\lean{BooleanAnalysis.ThresholdFunctions.affinePoly_singleton}
\leanok
\uses{BooleanAnalysis.ThresholdFunctions.affinePoly}
\difficulty{1}
Let $n\ge0$. Let $c\in\mathbb R$, $a\in\mathbb R^n$, and let $A_{c,a}(S)$ equal $c$ for $S=\varnothing$, $a_i$ for $S=\{i\}$, and zero for all larger $S$. For every $i\in[n]$, $A_{c,a}(\{i\})=a_i$.
\end{lemma}

\begin{lemma}[Degree bound for the affine coefficient family]
\lean{BooleanAnalysis.ThresholdFunctions.affinePoly_hasDegreeAtMost_one}
\leanok
\uses{BooleanAnalysis.ThresholdFunctions.MultilinearPolynomial.HasDegreeAtMost, BooleanAnalysis.ThresholdFunctions.affinePoly}
\difficulty{3}
Let $n\ge0$. Let $c\in\mathbb R$, $a\in\mathbb R^n$, and let $A_{c,a}(S)$ equal $c$ for $S=\varnothing$, $a_i$ for $S=\{i\}$, and zero for all larger $S$. This coefficient family has degree at most one: $A_{c,a}(S)=0$ whenever $|S|>1$.
\end{lemma}

\begin{lemma}[Affine representation of a linear threshold function]
\lean{BooleanAnalysis.ThresholdFunctions.IsLinearThreshold.exists_affine}
\leanok
\uses{BooleanAnalysis.BooleanFunc, BooleanAnalysis.ThresholdFunctions.IsLinearThreshold, BooleanAnalysis.ThresholdFunctions.MultilinearPolynomial.eval_of_degree_le_one, BooleanAnalysis.ThresholdFunctions.MultilinearPolynomial.threshold, BooleanAnalysis.ThresholdFunctions.thresholdSign, BooleanAnalysis.boolToSign}
\difficulty{3}
Let $n\ge0$ and give $\{0,1\}^n$ uniform probability; write $\mathbb E$ for expectation. Put $\chi_S(x)=\prod_{i\in S}(-1)^{x_i}$ and $\widehat h(S)=\mathbb E[h\chi_S]$ for $S\subseteq[n]$ and real cube functions $h$. Write $\operatorname{sgn}_+(t)=1$ for $t\ge0$ and $-1$ for $t<0$. If $f$ is the sign $\operatorname{sgn}_+(P)$ of a Walsh polynomial $P$ whose coefficients vanish on sets larger than one, then there exist real $c$ and $a_1,\ldots,a_n$ such that $f(x)=\operatorname{sgn}_+(c+\sum_i a_i(-1)^{x_i})$ for every $x$.
\end{lemma}

\begin{lemma}[An affine sign is a linear threshold function]
\lean{BooleanAnalysis.ThresholdFunctions.isLinearThreshold_affine}
\leanok
\uses{BooleanAnalysis.ThresholdFunctions.IsLinearThreshold, BooleanAnalysis.ThresholdFunctions.MultilinearPolynomial.eval_of_degree_le_one, BooleanAnalysis.ThresholdFunctions.MultilinearPolynomial.threshold, BooleanAnalysis.ThresholdFunctions.affinePoly, BooleanAnalysis.ThresholdFunctions.affinePoly_hasDegreeAtMost_one, BooleanAnalysis.ThresholdFunctions.thresholdSign, BooleanAnalysis.boolToSign}
\difficulty{3}
Let $n\ge0$ and give $\{0,1\}^n$ uniform probability; write $\mathbb E$ for expectation. Write $\operatorname{sgn}_+(t)=1$ for $t\ge0$ and $-1$ for $t<0$. For every real $c$ and $a\in\mathbb R^n$, the function $x\mapsto\operatorname{sgn}_+(c+\sum_i a_i(-1)^{x_i})$ is a linear threshold function, meaning the sign of a Walsh polynomial with no coefficient above degree one.
\end{lemma}

\begin{lemma}[Separating a fixed coordinate of an affine form]
\lean{BooleanAnalysis.ThresholdFunctions.affine_form_update}
\leanok
\uses{BooleanAnalysis.BoolCube, BooleanAnalysis.boolToSign}
\difficulty{3}
Let $n\ge0$, $c\in\mathbb R$, $a\in\mathbb R^n$, $i\in[n]$, $x\in\{0,1\}^n$, and $b\in\{0,1\}$. If $x_{i\leftarrow b}$ replaces bit $i$ by $b$, then $c+\sum_j a_j(-1)^{(x_{i\leftarrow b})_j}=(c+\sum_{j\ne i}a_j(-1)^{x_j})+a_i(-1)^b$.
\end{lemma}

\begin{lemma}[Every affine threshold derivative has a constant sign]
\lean{BooleanAnalysis.ThresholdFunctions.derivative_affine_threshold_sign}
\leanok
\uses{BooleanAnalysis.ThresholdFunctions.affine_form_update, BooleanAnalysis.ThresholdFunctions.thresholdSign, BooleanAnalysis.ThresholdFunctions.thresholdSign_mono, BooleanAnalysis.boolToSign, BooleanAnalysis.boolToSign_false, BooleanAnalysis.boolToSign_true, BooleanAnalysis.derivative}
\difficulty{4}
Let $n\ge0$ and give $\{0,1\}^n$ uniform probability; write $\mathbb E$ for expectation. Write $\operatorname{sgn}_+(t)=1$ for $t\ge0$ and $-1$ for $t<0$. For a real cube function $h$, write $D_i h(x)=\tfrac12(h(x_{i\leftarrow0})-h(x_{i\leftarrow1}))$, where $x_{i\leftarrow b}$ replaces bit $i$ by $b$. For any real $c$, $a\in\mathbb R^n$, and $i\in[n]$, put $f(x)=\operatorname{sgn}_+(c+\sum_j a_j(-1)^{x_j})$. Either $D_i f(x)\ge0$ for every $x$, or $D_i f(x)\le0$ for every $x$.
\end{lemma}

\begin{lemma}[Singleton Fourier energy of a sign function]
\lean{BooleanAnalysis.ThresholdFunctions.sum_singleton_fourier_sq_le_one}
\leanok
\uses{BooleanAnalysis.BooleanFunc, BooleanAnalysis.fourierCoeff, BooleanAnalysis.isPmOne, BooleanAnalysis.parseval_pm_one}
\difficulty{3}
Let $n\ge0$ and give $\{0,1\}^n$ uniform probability; write $\mathbb E$ for expectation. Put $\chi_S(x)=\prod_{i\in S}(-1)^{x_i}$ and $\widehat h(S)=\mathbb E[h\chi_S]$ for $S\subseteq[n]$ and real cube functions $h$. Every $f\colon\{0,1\}^n\to\{-1,1\}$ satisfies $\sum_{i=1}^n\widehat f(\{i\})^2\le1$.
\end{lemma}

\begin{theorem}[Total influence bound for linear thresholds]
\lean{BooleanAnalysis.ThresholdFunctions.ltf_totalInfluence_le_sqrt}
\leanok
\uses{BooleanAnalysis.BooleanFunc, BooleanAnalysis.ThresholdFunctions.IsLinearThreshold, BooleanAnalysis.ThresholdFunctions.IsLinearThreshold.exists_affine, BooleanAnalysis.ThresholdFunctions.derivative_affine_threshold_sign, BooleanAnalysis.ThresholdFunctions.influence_eq_abs_fourierCoeff_of_sign, BooleanAnalysis.ThresholdFunctions.isPmOne_thresholdSign, BooleanAnalysis.ThresholdFunctions.sum_singleton_fourier_sq_le_one, BooleanAnalysis.ThresholdFunctions.thresholdSign, BooleanAnalysis.boolToSign, BooleanAnalysis.fourierCoeff, BooleanAnalysis.isPmOne, BooleanAnalysis.totalInfluence}
\difficulty{4}
Let $n\ge0$ and give $\{0,1\}^n$ uniform probability; write $\mathbb E$ for expectation. Write $\operatorname{sgn}_+(t)=1$ for $t\ge0$ and $-1$ for $t<0$. Set $I_i(h)=\tfrac14\mathbb E[(h(x)-h(x^{\{i\}}))^2]$ and $I(h)=\sum_i I_i(h)$, where $x^{\{i\}}$ flips bit $i$. If $f(x)=\operatorname{sgn}_+(c+\sum_i a_i(-1)^{x_i})$ for some real $c,a_i$, then $I(f)\le\sqrt n$.
\end{theorem}
